\documentclass[11pt]{article}
\renewcommand{\labelitemi}{--}

% Self-contained: uses only environments from the LaTeX kernel
% (verbatim, tabular, itemize, enumerate). No extra packages needed.

\begin{document}

\title{A Practical Path into Low-Level Programming}
\author{Pointers / Compilers / C / Memory Allocators}
\date{November 2026}
\maketitle

\section*{Overview}

You already have the right setup: Fedora, gcc 16, gdb, and a small SDL3 project
called \texttt{Pixels in C}. That project is the seed for the final exercise.
The stages below go from raw pointers to a miniature video player built on
FFmpeg and SDL.

\bigskip

\begin{center}
\begin{tabular}{@{}p{0.28\textwidth} p{0.64\textwidth}@{}}
\hline
\textbf{Stage} & \textbf{Deliverable} \\
\hline
0. Toolchain      & Compiler, debugger, sanitizers, headers \\
1. Pointers \& C  & String library, dynamic vector, linked list \\
2. Compilers      & Multi-file build, understand objects and linking \\
3. Memory         & Bump allocator, free-list allocator, refcounted buffer \\
4. Linux systems  \texttt{cat}, \texttt{mmap} copy, \texttt{fork}/\texttt{pipe}, threads \\
5. Capstone       & Mini-ffplay: FFmpeg decode --> SDL3 render \\
\hline
\end{tabular}
\end{center}

\section{Stage 0 --- Toolchain}

\begin{verbatim}
sudo dnf install gcc make gdb strace ffmpeg-devel
\end{verbatim}

\texttt{ffmpeg-devel} supplies the headers needed to compile against the
libraries already installed. Compile every exercise like this:

\begin{verbatim}
gcc -std=c11 -Wall -Wextra -Wpedantic -g \
    -fsanitize=address,undefined -o test test.c
\end{verbatim}

AddressSanitizer catches use-after-free, buffer overflows, double frees and
leaks at runtime. It replaces most of what valgrind is used for; valgrind is
not installed here, but gcc's sanitizers are.

\section{Stage 1 --- Pointers and C itself}

The mental model that makes everything click:

\begin{itemize}
  \item All memory is \textbf{bytes with addresses}. A pointer is an address
        plus a type that tells the compiler how many bytes to read or write.
  \item An array is not a separate entity: \texttt{a[i]} is \texttt{*(a+i)}.
        Arrays decay to pointers when passed to functions.
  \item A C string is a \texttt{char} array ending in a zero byte (NUL).
        There is no string type.
  \item \texttt{sizeof(p)} is the size of the pointer (8), not the pointee.
        Use \texttt{sizeof(*p)} for the latter.
  \item \texttt{void *} is a generic address; C converts it to and from any
        object pointer implicitly (unlike C++).
\end{itemize}

Exercises, in order:

\begin{enumerate}
  \item \texttt{void swap(int *a, int *b)} and print addresses with
        \texttt{printf("\%p", (void*)p)}.
  \item Reimplement \texttt{strlen}, \texttt{strcpy}, \texttt{strcat},
        \texttt{strcmp} without using libc string functions.
  \item A dynamic array: \verb|struct Vec { void *data; size_t len, cap; }|
        with \verb|vec_push()| built on \texttt{realloc()}.
  \item A singly linked list with a \verb|list_free()| that leaks nothing
        (verify with ASan).
\end{enumerate}

Draw boxes and arrows on paper for exercises 3 and 4. It is faster than
staring at the code.

Recommended first book: C Programming: A Modern Approach by Knezek.
Alternative: K\&R chapters 1--5 if you like terse classics.

\section{Stage 2 --- What a compiler actually does}

Run each command on a tiny \texttt{hello.c} and inspect the output:

\bigskip

\begin{center}
\begin{tabular}{@{}p{0.38\textwidth} p{0.54\textwidth}@{}}
\hline
\textbf{Command} & \textbf{What it shows} \\
\hline
\texttt{gcc -E hello.c}   & Preprocessed source: macros and includes expanded \\
\texttt{gcc -S hello.c}   & Generated assembly \\
\texttt{gcc -c a.c -o a.o}& Object file: machine code plus symbol table \\
\texttt{nm a.o}           & Symbols defined and referenced \\
\texttt{objdump -d a.o}   & Disassembly \\
\texttt{readelf -h a.o}   & ELF headers \\
\texttt{ar rcs libmine.a a.o b.o} & Static library \\
\texttt{gcc -v ...}       & Full pipeline: cpp --> cc1 --> as --> ld \\
\hline
\end{tabular}
\end{center}

Then split \texttt{pilot.c} into \texttt{pilot.c} and \texttt{pilot.h}, add
include guards, and deliberately cause an undefined reference error and
a multiple definition error. Learning to read those two messages saves
hours later.

Use godbolt.org to compare \texttt{-O0} and \texttt{-O2} side by side. That is
where pointers stop being abstract: \texttt{a[i]} becomes a base register plus
an offset.

This is also the answer to ``how does FFmpeg work'': \texttt{libavcodec.so} is
a pile of object files exporting symbols such as \verb|avcodec_find_decoder|.
VLC links against it and calls those symbols. No magic.

\section{Stage 3 --- Memory model and allocators}

Internalize the process layout:

\begin{verbatim}
high addresses
 +----------------+
 | stack          |  grows down: locals, return addresses
 +----------------+
 | (gap)          |
 +----------------+
 | heap           |  grows up: malloc / free / realloc
 +----------------+
 | .data / .bss   |  globals, static variables
 +----------------+
 | .text          |  machine code, read-only
 +----------------+
low addresses
\end{verbatim}

See your own version with \texttt{cat /proc/self/maps}.

Key facts:

\begin{itemize}
  \item \texttt{malloc()} asks the OS for pages (\texttt{mmap}/\texttt{brk})
        and carves them into chunks. It is not itself a syscall.
  \item \textbf{Alignment}: allocations are padded to 8 or 16 bytes.
        \verb|(p + 15) & ~15| rounds a pointer up to 16.
  \item \textbf{Padding}: \verb|struct { char c; int x; }| is 8 bytes,
        not 5. Field order matters.
  \item \texttt{realloc()} may move your data; old pointers into that block
        become dangling.
  \item Refcounting is how FFmpeg avoids copying frames:
        \texttt{AVFrame->buf[]} holds \texttt{AVBufferRef}s.
\end{itemize}

Allocator exercises:

\begin{enumerate}
  \item \textbf{Bump allocator}: static \texttt{char arena[65536]}, a
        \verb|bump_alloc()| that advances an offset with alignment. No free.
  \item \textbf{Free-list allocator}: a small header before each block
        (\texttt{size}, \texttt{next}), implement \verb|my_free()|, then
        coalesce adjacent free blocks.
  \item \textbf{Refcounted buffer}: \verb|struct Buf { void *data; size_t size; int *refs; }|
        with \verb|buf_ref()| and \verb|buf_unref()| ---
        a miniature \texttt{AVBufferRef}.
  \item Write buggy versions on purpose and watch ASan catch them.
\end{enumerate}

For the ``why'': Computer Systems: A Programmer's Perspective (CSAPP),
chapters 2, 3, 5 and 6.

\section{Stage 4 --- The Linux systems layer}

Learn these calls and their man pages: \texttt{open/read/write/close},
\texttt{mmap}, \texttt{fork/exec/wait}, \texttt{pipe}, \verb|pthread_create|,
mutexes and atomics. Trace them with \texttt{strace -e trace=file ./app}.

Exercises: a 20-line \texttt{cat}; a \texttt{cp} implemented with \texttt{mmap};
a \texttt{fork()} plus \texttt{pipe()} pair; and a shared counter with and
without a mutex so you can see the race.

Book: Advanced Programming in the UNIX Environment (APUE).

\section{Stage 5 --- Capstone: your own mini-ffplay}

This ties everything together with what you already have:

\begin{enumerate}
  \item After installing \texttt{ffmpeg-devel}, find the header location:
        \texttt{rpm -ql ffmpeg-devel | grep avcodec.h}.
  \item Compile \verb|doc/examples/demux_decode.c| from the FFmpeg source
        against those headers.
  \item Instead of writing raw frames to a file, feed them into your
        \texttt{pilot.c} SDL3 texture.
  \item Now you have demux --> decode --> scale --> render: ffplay in
        miniature.
  \item Bonus: route frame buffers through your own arena/refcount allocator.
\end{enumerate}

Do not wait until you have ``finished'' all of this before touching FFmpeg.
Do stages 1--3, then read \verb|doc/examples/demux_decode.c| alongside
stage 4. Keep one small program you own and add one layer at a time --- which
is exactly what \texttt{Pixels in C} already is.

\section{Four-week rhythm}

\bigskip

\begin{center}
\begin{tabular}{@{}p{0.12\textwidth} p{0.80\textwidth}@{}}
\hline
\textbf{Week} & \textbf{Focus and deliverable} \\
\hline
1 & Pointers, strings, gdb, warnings --- string library plus \texttt{Vec} \\
2 & Preprocessor, assembly, objects, linking, Makefiles --- multi-file build \\
3 & Memory layout, ASan, custom allocators, padding --- bump and free-list allocators \\
4 & POSIX I/O, fork, threads, FFmpeg example --- mini-player with SDL3 \\
\hline
\end{tabular}
\end{center}

\end{document}
